\section{Reference Orbit Calculation}
\label{sec:ref-orbit-calc}

Reference-orbit construction can dominate the cost of
deep-zoom Mandelbrot rendering.  Perturbation theory
(\cref{sec:perturbation-concept}) allows every pixel to be rendered at
low precision, but it requires one orbit---the \emph{reference
orbit}---whose accuracy resolves the zoom depth. The usual fixed backends
compute that orbit at the full configured precision. An optional dynamic MT9
variant reduces working precision when independently certified sample
conversions permit it (\cref{subsec:ref-orbit-dynamic-precision}). At
extreme magnifications this requirement means iterating
\begin{equation}
  z_{n+1} = z_n^2 + c,\qquad z,c\in\mathbb{C},
  \label{eq:mandelbrot-recurrence}
\end{equation}
for potentially \emph{billions} of steps, with each step carried out at
\emph{thousands} of bits of precision.  The resulting orbit may contain
tens of billions of elements and occupy tens of gigabytes of memory.
Computing it efficiently, compressing it to fit in RAM, and reusing it
across multiple renders at different perturbation precisions are
therefore critical to the overall performance of the system.

The following sections describe how \code{RefOrbitCalc} orchestrates single-threaded and
multithreaded MPIR backends and the CUDA implementation,
together with storage and compression modes that trade memory footprint against
recomputation.  The CUDA implementation uses a fused NTT recurrence while
preserving the application contract and mathematical recurrence of the CPU
backends.

\subsection{High-level pipeline and dispatch}
\label{subsec:ref-orbit-pipeline}

A reference orbit is stored in a \code{PerturbationResults<IterType,T,PExtras>}
instance, where:

\begin{itemize}
  \item \code{IterType} is the iteration index type (\code{uint32\_t} or \code{uint64\_t}).
  \item \code{T} is the low-precision numeric type used for downstream perturbation math (e.g.\ \code{float},
        \code{double}, \code{HDRFloat<...>}).
  \item \code{PExtras} selects the storage format for per-iteration orbit data:
        uncompressed (\code{Disable}), lightweight waypoint compression
        (\code{SimpleCompression}), Imagina-compatible maximal compression
        (\code{MaxCompression}), or diagnostic bookkeeping (\code{Bad}).
\end{itemize}

The two precisions have different jobs: MPIR or \code{HpSharkFloat} resolves
the location of the reference at the zoom depth, while \code{T} stores the
coefficients consumed by nearby pixel orbits.  Computing a reference at high
precision does not imply that every stored sample has that same precision.

Orbit construction is initiated through \code{AddPerturbationReferencePoint()}, which:

\begin{enumerate}
  \item Picks an initial guess \((c_x,c_y)\) (center of the current view
        if unset).
  \item Chooses an algorithm from \code{m\_Perturbation\-Alg}.  In
        \code{Auto}, the estimated precision selects \code{STPeriodicity}
        below $640$ bits, \code{MTPeriodicity3} from $640$ through $6{,}783$ bits,
        \code{MTPeriodicity3PerturbMTHighMTMed3} from $6{,}784$ through
        $16{,}383$ bits, and dynamic MT9 from $16{,}384$
        through $32{,}767$ bits. At $32{,}768$ bits and above, \code{GPU} is
        selected when CUDA is available and the effective render algorithm
        supports GPU reference generation within its precision capacity;
        otherwise, dynamic MT9 is selected. The estimate includes the minimum precision
        margin, excludes reuse padding, and does not depend on hardware concurrency.
        Explicit selections bypass Auto. Explicit fixed MT9 disables dynamic
        precision even at depths where Auto selects it.
  \item Allocates a new \code{PerturbationResults} slot, initializes
        metadata and bounds, and runs the chosen orbit kernel until escape,
        periodicity detection, or the maximum iteration count is reached.
\end{enumerate}

To control memory pressure, \code{OptimizeMemory()} monitors process commit
usage and opportunistically drops cached orbits that are not of the currently
demanded variant type when the working set exceeds a configurable threshold.

\paragraph{Complete selection catalog.}
\Cref{tab:reference-backends} lists every \code{PerturbationAlg} value and its
actual dispatch. A reuse selection controls reference generation and subsequent
medium-precision reference construction, rather than selecting the per-pixel
render algorithm. Reuse storage is described in \cref{subsec:ref-orbit-reuse}.
The legacy menu label ``MT2'' counts the two square workers; the implementation
name MT3 also counts the calling thread. Explicit selections bypass Auto.
Every selection name in the catalog is accepted by
\code{-{}-perturbation-alg}, including
\code{MTPeriodicity9\allowbreak Karatsuba\allowbreak DynamicPrecision}.
The CLI parser and help share one catalog, whose compile-time checks require
one unique, nonempty name for every algorithm enum value.

\begin{table}[!htbp]
\centering
\small
\caption[Reference-orbit selections]{Complete reference-orbit selection catalog.
  The CPU periodicity test is a termination heuristic, not an accuracy certificate.}
\label{tab:reference-backends}
\begin{tabularx}{\textwidth}{>{\raggedright\arraybackslash}p{0.48\textwidth}X}
\hline
\textbf{Selection name} & \textbf{Actual behavior} \\
\hline
\code{ST} & Factored MPIR recurrence, no periodicity. \\
\code{MT} & MT3 expanded recurrence, no periodicity. \\
\code{STPeriodicity} & ST with periodicity. \\
\code{MTPeriodicity3} & Two square workers plus caller, with periodicity. \\
\code{MTPeriodicity3\-PerturbMTHigh\-STMed} & ST authoritative generation,
  \code{SaveForReuse3}; single-threaded medium reference. \\
\code{MTPeriodicity3\-PerturbMTHigh\-MTMed1} & MT3 authoritative generation,
  uncompressed \code{SaveForReuse1}; caller stores reuse entries. \\
\code{MTPeriodicity3\-PerturbMTHigh\-MTMed2} & MT3 authoritative generation,
  uncompressed \code{SaveForReuse2}; additional reuse-storage worker. \\
\code{MTPeriodicity3\-PerturbMTHigh\-MTMed3} & MT3 authoritative generation,
  compressed \code{SaveForReuse3}; additional reuse-storage worker. \\
\code{MTPeriodicity3\-PerturbMTHigh\-MTMed4} & Broken reuse-v4 selection;
  not a supported reference backend. \\
\code{MTPeriodicity5} & Disabled MT5 entry point delegates to periodic MT3. \\
\code{GPU} & Fused CUDA Ref2/v2 recurrence; explicit or automatic selection. \\
\code{Auto} & Bit thresholds above; compatible GPU generation from
  $32{,}768$ bits, with dynamic MT9 as the CPU alternative. \\
\code{MTPeriodicity9\-Karatsuba} & Eight workers plus caller; exact split products,
  full-precision recurrence additions and subtractions, with periodicity. \\
\code{MTPeriodicity9\-Karatsuba\-DynamicPrecision} & Explicit GUI/CLI or Auto selection;
  variable working precision with sample certificates and fixed-MT9 restart.
  Fixed MT9 remains separately selectable. \\
\hline
\end{tabularx}
\end{table}

GPU eligibility in Auto is derived from the typed render catalog and the existing
CUDA runtime state. Supported HDR reference consumers include CPU renderers when
CUDA is available; reference generation and per-pixel rendering remain separate
choices. Scaled HDR rendering requires diagnostic reference bookkeeping, and
compressed double-float rendering lacks the required plain-to-compressed
materialization before conversion. Those formats use dynamic MT9 in Auto. Requested
precision above $16{,}777{,}216$ bits also uses dynamic MT9 because the largest GPU
bucket contains $524{,}288$ 32-bit limbs. Auto therefore avoids selecting a GPU
bucket that would cap the requested precision. The lower thresholds round the
former $10^{150}$ and $10^{2000}$ zoom rules to 64-bit limb boundaries.

\subsection{Single-threaded authoritative orbit}
\label{subsec:ref-orbit-st}

The single-threaded path is
\code{Add\-Perturbation\-Reference\-Point\-ST}.
It computes the application-authoritative orbit using finite, explicitly selected
precision in
MPIR/GMP~\cite{MPIR_Library} \code{mpf\_t} for the recurrence, while
simultaneously emitting a low-precision shadow copy
\((\hat{x}_n,\hat{y}_n)\in T^2\) for downstream work (compression,
bailout checks, periodicity tests).

\paragraph{State and initialization.}

Given a selected reference parameter \(c = c_x + i c_y\), the implementation:

\begin{itemize}
  \item Initializes \code{mpf\_t} temporaries for \(x,y\), scaled operands, and scratch products.
  \item Sets the loop-entry iterate to \((x_1,y_1) = (c_x,c_y)\), which
        corresponds to folding in the first iteration of the standard
        recurrence \(z_{n+1}=z_n^2+c\) with \(z_0=0\).
  \item Computes low-precision cast values \(\hat{c}_x,\hat{c}_y \in T\) either via \code{mpf\_get\_d} for
        native float/double, or via a mantissa/exponent extraction for extended formats.
\end{itemize}

\paragraph{Recurrence.}
Writing \(z_n = x_n + i y_n\), the actual single-threaded schedule evaluates
\begin{align}
  y_{n+1} &= (2x_n)y_n + c_y, \label{eq:st-recurrence-imag} \\
  x_{n+1} &= (x_n+y_n)(x_n-y_n) + c_x. \label{eq:st-recurrence-real}
\end{align}
The identity $(x_n+y_n)(x_n-y_n)=x_n^2-y_n^2$ gives the same complex
polynomial as \eqref{eq:mandelbrot-recurrence}. It removes one large product
from the serial instruction sequence: two ordinary MPIR \code{mpf\_mul}
calls replace two squares and a cross product. Forming the sum, difference,
and doubled operand still costs linear work in the limb count.

One calling thread owns every operation; \code{STPeriodicity} adds derivative
tracking to that same thread, not another worker. Its ordered loop body is:
\begin{enumerate}
  \item Form $2x_n$ with \code{mpf\_mul\_2exp}. Extract the old-state shadow
        $(\hat{x}_n,\hat{y}_n)$, emit its sample or compression input, and,
        when requested, capture the authoritative reuse entry and test periodicity.
  \item Multiply $(2x_n)y_n$ and add $c_y$ into a scratch destination.
        The live $y$ buffer still holds $y_n$.
  \item Form $x_n+y_n$ and $x_n-y_n$ in separate temporaries.
  \item Multiply those factors into the live $x$ destination and add $c_x$.
  \item Swap the live $y$ buffer with the completed imaginary scratch buffer.
        The swap changes buffer ownership without copying all its limbs.
  \item Apply the low-precision bailout test and the loop's stopping checks.
\end{enumerate}
An early periodicity return can terminate the body before the arithmetic update.
Otherwise both coordinates use the same old state, even though $x$ is overwritten
before the new $y$ is installed. Persistent MPIR temporaries are initialized
before the loop; result storage can grow separately. On exit, MPIR objects
allocated under a thread-local allocator are destroyed before its shutdown.
\Cref{fig:cpu-reference-dependencies} compares this serial schedule with MT3.

\paragraph{Bailout.}

The bailout threshold is evaluated in low precision using

\begin{equation}
  \lVert\hat{z}_n + \hat{c}\rVert^2 = (\hat{x}_n + \hat{c}_x)^2 + (\hat{y}_n + \hat{c}_y)^2 > 256,
  \label{eq:bailout-test}
\end{equation}

which matches the code's use of \code{TwoFiftySix} and avoids a high-precision
norm each step.  Here \(\hat{x}_n\) and \(\hat{y}_n\) are the low-precision
shadows of the orbit state captured \emph{before} the high-precision recurrence
update (i.e.\ they represent the current \(z_n\), not the just-computed
\(z_{n+1}\)).  The sum \(\hat{z}_n + \hat{c}\) is therefore a cheap proxy that
avoids re-extracting a low-precision snapshot after the update.
The proxy is not the next iterate: for $c=1$ and $z_n=2$,
$z_n+c=3$, whereas $z_{n+1}=5$.  Its threshold is an application termination
convention for reference generation, distinct from the mathematical escape
radius and the final pixel bailout.  Consequently, reference and pixel loops
need not terminate at the same index.

\subsection{View-dependent near-return checking}
\label{subsec:ref-orbit-periodicity}

An orbit returning close to the critical point can indicate a nearby periodic
center.  Its closeness must be interpreted relative to parameter sensitivity:
a displacement $\delta c$ changes $z_n$ to first order by $d_n\delta c$.
The CPU and GPU reference backends use that observation for the same
radius-based early-termination test. The following derivation describes their
shared mathematics; \cref{sec:hp-periodicity} describes GPU execution.
The implementations track the complex derivative

\(\frac{\partial z_n}{\partial c}\) in low precision:
\begin{equation}
  d_{n+1} = 2 z_n d_n + 1,\qquad d_0 = 0.
  \label{eq:dzdc-def}
\end{equation}
Mathematically \(d_0 = \partial z_0/\partial c = 0\), since \(z_0=0\).
The chain rule gives
$\partial(z_n^2+c)/\partial c=2z_n(\partial z_n/\partial c)+1$;
the final $+1$ is the direct dependence on the parameter.
Because ordinary reference generation begins its loop at \(z_1=c\) (folding in the
first iteration), the derivative's real component is initialized to~$1$,
which corresponds to \(d_1=1\) in the recurrence above.

Writing \(d_n = d_{x,n} + i d_{y,n}\) and expanding into real and imaginary parts yields:

\begin{align}
  d_{x,n+1} &= 2(x_n d_{x,n} - y_n d_{y,n}) + 1, \label{eq:st-dzdc-real} \\
  d_{y,n+1} &= 2(x_n d_{y,n} + y_n d_{x,n}). \label{eq:st-dzdc-imag}
\end{align}

\subsubsection{Screen coverage and parameter distance}
\label{subsec:ref-orbit-view-bounds}

Let the image cover the parameter rectangle
$[x_{\min},x_{\max}]\times[y_{\min},y_{\max}]$, with reference parameter
$c_\star=c_{\star,x}+ic_{\star,y}$. Bounds relative to that reference are
\begin{align}
 R_x&=\max(|x_{\min}-c_{\star,x}|,|x_{\max}-c_{\star,x}|),\\
 R_y&=\max(|y_{\min}-c_{\star,y}|,|y_{\max}-c_{\star,y}|).
 \label{eq:periodicity-screen-radii}
\end{align}
Every sampled pixel then satisfies
$|\Delta c_x|\le R_x$, $|\Delta c_y|\le R_y$, and
$|\Delta c|\le\rho=\sqrt{R_x^2+R_y^2}$. The rectangle includes all sample
positions, so these bounds remain conservative when the last sampled position
precedes an edge. An off-center or reused reference requires its actual
coordinates in these expressions.

For output dimensions $W\times H$ and $s$ samples per pixel dimension, the
coordinate steps are
\begin{equation}
 h_x=\frac{x_{\max}-x_{\min}}{sW},\qquad
 h_y=\frac{y_{\max}-y_{\min}}{sH}.
 \label{eq:periodicity-pixel-spacing}
\end{equation}
The covered widths are $sWh_x$ and $sHh_y$. Increasing resolution or
supersampling within the same rectangle refines sample spacing without
increasing its parameter-distance bound. For a centered reference and an
aspect-fitted view with square pixels, $R_x/R_y=W/H$.

\subsubsection{Linear displacement and the factor two}
\label{subsec:ref-orbit-linear-radius}

Write $d_n=a+ib$ and $\Delta c=u+iv$. Complex multiplication gives
$d_n\Delta c=(au-bv)+i(bu+av)$, so the rectangular bounds imply
\begin{equation}
 \|d_n\Delta c\|_\infty\le B_n
 =\max(|a|R_x+|b|R_y,\;|b|R_x+|a|R_y).
 \label{eq:periodicity-rectangular-bound}
\end{equation}
Each component bound is attained at a corner of the symmetric offset
rectangle. If both $R_x$ and $R_y$ are at most $R$, then
\begin{equation}
 B_n\le R(|a|+|b|)\le2R\|d_n\|_\infty.
 \label{eq:periodicity-factor-two}
\end{equation}
The factor two therefore has a direct norm-based interpretation for a
covering square in parameter space. A reference state small relative to
that displacement scale suggests a nearby critical return. The scalar
inequality uses an enclosing square for the linear image; it does not solve
$\Delta c=-z_n/d_n$ or establish that such an offset lies in the view.

The current CPU \code{GetMaxRadius()} and GPU \code{RadiusY} both store
the vertical half-height $(y_{\max}-y_{\min})/2$. The initializer does not
use the horizontal extents or reference offset for that scalar. Its name
therefore does not establish a maximum distance over the complete image.
For a centered $16{:}9$ view and $d_n=1+i$,
\eqref{eq:periodicity-rectangular-bound} gives $B_n=25R_y/9$, whereas
the implemented scale is $2R_y$. Equation~\eqref{eq:periodicity-factor-two}
describes a covering-square bound only when its radius actually covers
both axes. The existing implementation retains its vertical-radius policy.

\subsubsection{Nonlinear displacement and reduced arithmetic}
\label{subsec:ref-orbit-near-return-limits}

The screen bounds constrain $\Delta c$, not the evolved orbit displacement.
For an exact reference, the complete perturbation identity is
\begin{equation}
 \Delta z_{n+1}=2z_n^\star\Delta z_n+(\Delta z_n)^2+\Delta c.
 \label{eq:periodicity-full-displacement}
\end{equation}
Consequently, a uniform disk bound in exact arithmetic can be propagated as
\begin{equation}
 D_0=0,\qquad
 D_{n+1}=2|z_n^\star|D_n+D_n^2+\rho,
 \qquad |\Delta z_n|\le D_n.
 \label{eq:periodicity-displacement-disk}
\end{equation}
The triangle inequality proves the induction. Such a bound can be much wider
than the linear estimate; a numerical implementation would additionally
require outward rounding and reference-error allowances. The existing
near-return test does not propagate this disk or certify a nonlinear zero.

The geometric bounds are independent of the numeric type. Their computed
values, orbit snapshots, derivative recurrence, and final comparison incur
type-dependent errors. \code{HDRFloat<float>} and \code{HDRFloat<double>}
retain 24 and 53 significand bits, respectively; the separate exponent extends
range without supplying more significand bits. Derivative cancellation and
accumulation prevent a universal relative-error allowance based on machine
epsilon alone. A comparison whose true slack
$2R\|d_n\|_\infty-\|z_n\|_\infty$ is close to zero can therefore change
with the reduced type. More precision improves its evaluation but does not
replace the geometric or nonlinear arguments.

\subsubsection{Implemented comparison and CPU execution}

The shared reduced-precision comparison forms

\begin{equation}
  n_2 = \max(|\hat{x}_n|,|\hat{y}_n|),\qquad
  r_0 = \max(|\hat d_{x,n}|,|\hat d_{y,n}|),
  \label{eq:periodicity-norms}
\end{equation}

and define a detection threshold

\begin{equation}
  n_3 = 2\,R_{\max}\,r_0,
  \label{eq:periodicity-threshold}
\end{equation}

where $R_{\max}$ denotes the stored vertical half-height described above,
and hats denote the reduced values actually compared. The completed threshold
is normalized before the HDR exponent-first comparison. If \(n_2 < n_3\),
the CPU orbit is marked as \emph{maybe periodic} and the
reference loop terminates early (unless periodicity detection is disabled).
Otherwise, \((d_{x,n},d_{y,n})\) is advanced using the update above.
Single-threaded generation checks its current snapshot before updating the
high-precision state. MT3 and MT9 retain their delayed sample schedules
(\cref{subsec:ref-orbit-mt3,subsec:ref-orbit-mt9}); their arithmetic and
publication timings need not match GPU execution bit for bit.

Early reference termination does not terminate every pixel's iteration.
At the end of an available reference prefix, rebasing represents the current
pixel state relative to $z_0^\star=0$ and continues its own iteration count
(\cref{sec:perturb-only-deltas,sec:perturb-only-reconstruct}). That algebraic
change of representation does not require an exactly periodic reference.
Finite-precision rendering and compression retain their existing numerical
contracts. The independent dynamic-MT9 sample certificate in
\cref{subsec:ref-orbit-dynamic-certificate} certifies conversions rather
than this near-return decision.

For $c_\star=-1$, the complete calculation remains real:
\begin{center}
\begin{tabular}{c|rrrrr}
$n$ & 0 & 1 & 2 & 3 & 4 \\
\hline
$z_n$ & 0 & $-1$ & 0 & $-1$ & 0 \\
$d_n$ & 0 & 1 & $-1$ & 1 & $-1$
\end{tabular}
\end{center}
For example, $d_2=2(-1)(1)+1=-1$.  At $n=2$ any positive
$R_{\max}$ makes $0<2R_{\max}|d_2|$ true.  The exact orbit here has
period two, but applying the same reduced-precision inequality at a nearby
parameter does not prove an exact return, identify the least period, or
certify every pixel's accuracy.  The stored ``maybe periodic'' status retains
that distinction.

\subsection{Expanded versus factored evaluation of the reference orbit}
\label{subsec:ref-orbit-expanded-factored}

The high-precision reference orbit evaluates the quadratic map

\[
z_{n+1} = z_n^2 + c,
\qquad
z_n = x_n + i y_n,
\qquad
c = c_x + i c_y,
\]

with the finite MPIR or fixed-width \code{HpSharkFloat} precision selected by
the active backend.

Writing the iteration in real and imaginary components gives the
\emph{expanded form}:

\begin{align}
x_{n+1} &= x_n^2 - y_n^2 + c_x, \label{eq:ref-expanded-real} \\
y_{n+1} &= 2 x_n y_n + c_y. \label{eq:ref-expanded-imag}
\end{align}

The expanded form follows directly from the algebraic definition of the polynomial.
Each term is evaluated explicitly, requiring two full-precision squares and
one full-precision multiplication per iteration.

The same quadratic map can be evaluated in a mathematically equivalent
\emph{factored form}. In particular, the real component may be written as

\begin{align}
x_{n+1}
&= (x_n - y_n)(x_n + y_n) + c_x, \label{eq:ref-factored-real}
\end{align}
while the imaginary component remains
\begin{align}
y_{n+1} &= 2 x_n y_n + c_y. \label{eq:ref-factored-imag}
\end{align}

In exact arithmetic, \cref{eq:ref-factored-real,eq:ref-factored-imag} are
identical to \cref{eq:ref-expanded-real,eq:ref-expanded-imag}. The difference
lies solely in how the products are formed.

From an implementation perspective, the expanded form evaluates
$x_n^2$ and $y_n^2$ independently and makes the subtraction
$x_n^2 - y_n^2$ explicit. This closely mirrors the mathematical definition
of the Mandelbrot polynomial and exercises the full squaring and
multiplication paths of the high-precision arithmetic.

For serial evaluation, factoring shortens the straight-line arithmetic sequence
from three large products to two. The saving is one product, not a shorter
Mandelbrot orbit or a reduction in required precision. Its value depends on
the relative costs of a specialized square, a general multiplication, and the
extra sum/difference formation. The independent expanded squares instead
expose parallel tasks, which explains why the CPU backends use different
schedules. Factoring also changes where cancellation and rounding occur;
exact polynomial equality is not an equality of finite-precision instruction
sequences. \Cref{subsec:ref-orbit-mt-tradeoff} accounts for these distinct costs.

For $x_n=3/4$, $y_n=1/2$, and $c=0$, the expanded real update is
$9/16-1/4=5/16$.  The factored update is
$(3/4-1/2)(3/4+1/2)=(1/4)(5/4)=5/16$; both give imaginary part
$2(3/4)(1/2)=3/4$.  Equality of these exact results does not require the
rounded intermediate products to agree for arbitrary finite-precision inputs.

The authoritative single-threaded CPU path uses the factored schedule through
\code{mpf\_mul}; \code{MT3} uses independent $x_n^2$ and
$y_n^2$ workers and a main-thread cross product. MT9 instead reconstructs
each of those three products from exact half-width integer products
(\cref{subsec:ref-orbit-mt9}). The CUDA engine
implements the same component equations through the factored nonlinear products
\begin{align}
R_n &= (x_n+y_n)(x_n-y_n)=x_n^2-y_n^2, \\
I_n &= 2x_ny_n,
\end{align}
which it constructs directly in the NTT domain.  It then incorporates $c_x$ and
$c_y$ at their exact bit offsets during signed unpack, before positional carry
propagation.  The CUDA engine therefore never materializes three normalized
product values or invokes a standalone high-precision add/subtract pipeline.

Fusion does not eliminate the final base-$2^{32}$ dependency.  The CUDA engine forms every
nonlinear and linear contribution first, then resolves the completed real and
imaginary components together in one packed decoupled-look-back carry phase.
Each component remains an independent logical carry chain within that phase;
\cref{subsec:fused-ntt-reference-step,subsec:fused-ntt-periodicity} describes
the construction, operation schedule, carry mathematics, and periodicity test.
